\chapter{Orthogonal motives of rank two and quadratic forms}

This chapter transcribes \S8 (``Orthogonal motives of rank $2$'') and \S9
(``Quadratic forms'') of Ancona's paper. We first state the project's main
technical result, \cref{thm:mainthm}: for a rank-two orthogonal motive in mixed
characteristic, positivity of the Betti pairing forces positivity of the
pairing on algebraic classes over the residue field. This is the statement that
the proof of the standard conjecture of Hodge type (\cref{thm:scht}) reduces to;
that reduction is carried out in the chapter on the standard conjecture of Hodge
type and is \emph{not} repeated here. The proof of \cref{thm:mainthm} given here
assembles the \(p\)-adic period computations of the later chapters; those
computations themselves are established there and only invoked here.

We then recall the classical theory of rank-two \(\Q\)-quadratic forms --- the
Hilbert symbol, the local invariants, the product formula --- and combine them
into a sign-comparison corollary (\cref{cor:signature-top}). The chapter
culminates in \cref{prop:padic-reduction}, which reduces \cref{thm:mainthm} to a
single comparison of the two quadratic spaces over \(\Q_p\).

Throughout, \(K\) is a \(p\)-adic field with ring of integers \(W\) and residue
field \(k\), and \(\Q_\nu\) denotes the completion of \(\Q\) at the place
\(\nu\).

\section{Orthogonal motives of rank $2$}

\subsection*{The main technical result}

\begin{definition}[Setup and hypotheses of the main theorem]
  \label{def:mainthm-setup}
  \lean{MR4199442StandardConjecturesForAbelianFourfolds.MainTheoremSetup}
  \uses{def:chow-motives, def:num-motives, def:realizations, def:pairing-mot, def:filtered-phi-modules}
  % SOURCE: [Ancona], §8, Theorem (mainthm) (read from references/source/fourfolds_2020.tex)
  % SOURCE QUOTE: "Let $K$ be a $p$-adic field, $W$  its ring of integers and $k$  its residue field. Let us fix an embedding $\sigma : K \hookrightarrow \C$.
  % Let \[M \in \CHM(W)_\Q\] be a motive in mixed characteristic and consider the motives induced by pullbacks $M_{\vert\C} \in \CHM(\C)_\Q$ and $M_{\vert{k}} \in \CHM(k)_\Q$. Consider $V_B$ and $V_Z$ the two $\Q$-vector spaces defined as
  % \[V_B = R_B(M_{\vert\C} ) \hspace{0.5cm}
  % \textrm{and} \hspace{0.5cm}
  % V_Z = \Hom_{\NUM(k)_\Q}(\one, M_{\vert{k}}).\]
  % Let $q$ be a quadratic form on $M$, by which we mean a morphism in $ \CHM(W)_\Q$  of the form
  % $q: \Sym^2 M  \longrightarrow \one.$
  %  Consider the two $\Q$-quadratic forms induced by $q$ on $V_B$ and $V_Z$ respectively,
  % \[q_B=R_B(q)  \hspace{0.5cm} \textrm{and}  \hspace{0.5cm} q_Z(\cdot)= (q_{\vert{k}} \circ \Sym^2(\cdot)).\]
  % Suppose that the following holds:
  % \begin{enumerate}
  % \item\label{assumptiondim2} The two $\Q$-vector spaces $V_B$  and $V_Z$ are of dimension $2$.
  % \item The pairing $q_B$ on $V_B$ is a polarization of Hodge structures.
  % \end{enumerate}
  % Then the quadratic form $q_Z$ is positive definite."
  \textit{Source: Ancona, arXiv:1806.03216, §8 (hypotheses of the main theorem, bundled).}
  We bundle once and for all the data and hypotheses shared by the main theorem
  and every statement that ``keeps its notation''. Let \(K\) be a \(p\)-adic
  field, \(W\) its ring of integers and \(k\) its residue field, and fix an
  embedding \(\sigma : K \hookrightarrow \C\). Let
  \[M \in \CHM(W)_\Q\]
  be a motive in mixed characteristic, with induced motives
  \(M_{\vert\C} \in \CHM(\C)_\Q\) and \(M_{\vert{k}} \in \CHM(k)_\Q\) obtained by
  pullback. Consider the two \(\Q\)-vector spaces
  \[V_B = R_B(M_{\vert\C}) \qquad \text{and} \qquad
    V_Z = \Hom_{\NUM(k)_\Q}(\one, M_{\vert{k}}).\]
  Let \(q\) be a \emph{quadratic form on \(M\)}, that is, a morphism
  \(q : \Sym^2 M \to \one\) in \(\CHM(W)_\Q\), and let
  \[q_B = R_B(q) \qquad \text{and} \qquad
    q_Z(\cdot) = q_{\vert{k}} \circ \Sym^2(\cdot)\]
  be the induced \(\Q\)-quadratic forms on \(V_B\) and \(V_Z\) respectively. We
  assume throughout that
  \begin{enumerate}
    \item the \(\Q\)-vector spaces \(V_B\) and \(V_Z\) are both of dimension
      \(2\); and
    \item the pairing \(q_B\) on \(V_B\) is a polarization of Hodge structures.
  \end{enumerate}
\end{definition}

\begin{theorem}[Main theorem on rank-two orthogonal motives]
  \label{thm:mainthm}
  \lean{MR4199442StandardConjecturesForAbelianFourfolds.mainTheorem}
  \uses{def:mainthm-setup, prop:padic-reduction, prop:ordinary-case, ass:i-positive, prop:endo-field, prop:iso-norm-unram, cor:period-unram, lem:not-norm-unram, prop:iso-norm-ram, cor:period-ram, lem:not-norm-ram, cor:period-wild, lem:not-norm-wild, cor:period-wild2, lem:not-norm-wild2}
  % SOURCE: [Ancona], §8, Theorem (mainthm) (read from references/source/fourfolds_2020.tex)
  % SOURCE QUOTE: "Let $K$ be a $p$-adic field, $W$  its ring of integers and $k$  its residue field. Let us fix an embedding $\sigma : K \hookrightarrow \C$.
  % Let \[M \in \CHM(W)_\Q\] be a motive in mixed characteristic and consider the motives induced by pullbacks $M_{\vert\C} \in \CHM(\C)_\Q$ and $M_{\vert{k}} \in \CHM(k)_\Q$. Consider $V_B$ and $V_Z$ the two $\Q$-vector spaces defined as
  % \[V_B = R_B(M_{\vert\C} ) \hspace{0.5cm}
  % \textrm{and} \hspace{0.5cm}
  % V_Z = \Hom_{\NUM(k)_\Q}(\one, M_{\vert{k}}).\]
  % Let $q$ be a quadratic form on $M$, by which we mean a morphism in $ \CHM(W)_\Q$  of the form
  % $q: \Sym^2 M  \longrightarrow \one.$
  %  Consider the two $\Q$-quadratic forms induced by $q$ on $V_B$ and $V_Z$ respectively,
  % \[q_B=R_B(q)  \hspace{0.5cm} \textrm{and}  \hspace{0.5cm} q_Z(\cdot)= (q_{\vert{k}} \circ \Sym^2(\cdot)).\]
  % Suppose that the following holds:
  % \begin{enumerate}
  % \item\label{assumptiondim2} The two $\Q$-vector spaces $V_B$  and $V_Z$ are of dimension $2$.
  % \item The pairing $q_B$ on $V_B$ is a polarization of Hodge structures.
  % \end{enumerate}
  % Then the quadratic form $q_Z$ is positive definite."
  \textit{Source: Ancona, arXiv:1806.03216, §8.}
  Let \(K\) be a \(p\)-adic field, \(W\) its ring of integers and \(k\) its
  residue field, and fix an embedding \(\sigma : K \hookrightarrow \C\). Let
  \[M \in \CHM(W)_\Q\]
  be a motive in mixed characteristic, with induced motives
  \(M_{\vert\C} \in \CHM(\C)_\Q\) and \(M_{\vert{k}} \in \CHM(k)_\Q\) obtained by
  pullback. Consider the two \(\Q\)-vector spaces
  \[V_B = R_B(M_{\vert\C}) \qquad \text{and} \qquad
    V_Z = \Hom_{\NUM(k)_\Q}(\one, M_{\vert{k}}).\]
  Let \(q\) be a \emph{quadratic form on \(M\)}, that is, a morphism
  \(q : \Sym^2 M \to \one\) in \(\CHM(W)_\Q\), and let
  \[q_B = R_B(q) \qquad \text{and} \qquad
    q_Z(\cdot) = q_{\vert{k}} \circ \Sym^2(\cdot)\]
  be the induced \(\Q\)-quadratic forms on \(V_B\) and \(V_Z\) respectively.
  Suppose that
  \begin{enumerate}
    \item the \(\Q\)-vector spaces \(V_B\) and \(V_Z\) are both of dimension
      \(2\); and
    \item the pairing \(q_B\) on \(V_B\) is a polarization of Hodge structures.
  \end{enumerate}
  Then the quadratic form \(q_Z\) is positive definite.
\end{theorem}

% SOURCE QUOTE PROOF: "Thank to Proposition \ref{ordinarycase}, we can work under Assumption \ref{ipos}, and fix the positive integer $i$ as in the assumption as well as the field $F$ from Proposition \ref{propendo}.
% Let us first suppose that  $F$ is unramified. We can then combine Proposition \ref{warmbrinon} and Corollary \ref{brinon} to conclude that the quadratic forms  $q_{B,p}$ and $q_{Z,p}$ are isomorphic if and only if $1/p^i$ is a norm. Because of Lemma \ref{notnorm}, this is the case   if and only if $i$ is even.
% Suppose now that $F$ is ramified but it is not $\mathbb{Q}_2(\sqrt{3})$ nor $\mathbb{Q}_2(\sqrt{-1})$. Then, following  Remark \ref{remtame}, we can combine Proposition \ref{warmbrinonram} and Corollary \ref{brinonram} to conclude that the quadratic forms $q_{B,p}$ and $q_{Z,p}$ are isomorphic if and only if    $ a^i$ is a norm. Because of Lemma \ref{notnormram}, this is the case   if and only if $i$ is even.
% If $F$ is  $\mathbb{Q}_2(\sqrt{-1})$ we can  combine Proposition \ref{warmbrinonram} and Corollary \ref{brinonwild} to conclude that the quadratic forms $q_{B,p}$ and $q_{Z,p}$ are isomorphic if and only if    $  (\alpha  \cdot \sigma(\alpha) )^i$ is a norm. Because of Lemma \ref{notnormwild}, this is the case   if and only if $i$ is even.
% Finally, if $F$ is  $\mathbb{Q}_2(\sqrt{3})$ we can  combine Proposition \ref{warmbrinonram} and Corollary \ref{brinonwild2} to conclude that the quadratic forms  $q_{B,p}$ and $q_{Z,p}$ are isomorphic if and only if    $  (\alpha  \cdot \sigma(\alpha) )^i$ is a norm. Because of Lemma \ref{notnormwild2}, this is the case   if and only if $i$ is even.
% In conclusion, we showed for that, for any possible $F$, the quadratic forms $q_{B,p}$ and $q_{Z,p}$ are isomorphic if and only if $i$ is even. This shows Theorem \ref{mainthm}  via Proposition \ref{padicreduction}."
\begin{proof}
  \uses{def:mainthm-setup, prop:padic-reduction, prop:ordinary-case,
    ass:i-positive, prop:endo-field, prop:iso-norm-unram, cor:period-unram,
    lem:not-norm-unram, prop:iso-norm-ram, cor:period-ram, lem:not-norm-ram,
    cor:period-wild, lem:not-norm-wild, cor:period-wild2, lem:not-norm-wild2}
  By \cref{prop:padic-reduction} it suffices to compare the two rank-two
  \(\Q_p\)-quadratic spaces \((V_{B,p}, q_{B,p})\) and \((V_{Z,p}, q_{Z,p})\):
  we must show that \(q_{B,p}\) and \(q_{Z,p}\) are isomorphic if and only if
  \(i\) is even, where \(i\) is the Hodge weight of \cref{prop:padic-reduction}.

  By \cref{prop:ordinary-case} we may assume that \(V_B\) is not of Hodge type
  \((0,0)\); that is, we may work under \cref{ass:i-positive}, fixing the
  positive integer \(i\) of that assumption and the field \(F\) furnished by
  \cref{prop:endo-field}. We now case-split on \(F\).

  \emph{\(F\) unramified.} Combining \cref{prop:iso-norm-unram} with the period
  computation of \cref{cor:period-unram}, the forms \(q_{B,p}\) and \(q_{Z,p}\)
  are isomorphic if and only if \(1/p^i\) is a norm; by
  \cref{lem:not-norm-unram} this holds if and only if \(i\) is even.

  \emph{\(F\) ramified, and \(F \neq \Q_2(\sqrt{3}), \Q_2(\sqrt{-1})\).}
  Combining \cref{prop:iso-norm-ram} with \cref{cor:period-ram}, the forms are
  isomorphic if and only if \(a^i\) is a norm; by \cref{lem:not-norm-ram} this
  holds if and only if \(i\) is even.

  \emph{\(F = \Q_2(\sqrt{-1})\).} Combining \cref{prop:iso-norm-ram} with
  \cref{cor:period-wild}, the forms are isomorphic if and only if
  \((\alpha \cdot \sigma(\alpha))^i\) is a norm; by \cref{lem:not-norm-wild}
  this holds if and only if \(i\) is even.

  \emph{\(F = \Q_2(\sqrt{3})\).} Combining \cref{prop:iso-norm-ram} with
  \cref{cor:period-wild2}, the forms are isomorphic if and only if
  \((\alpha \cdot \sigma(\alpha))^i\) is a norm; by \cref{lem:not-norm-wild2}
  this holds if and only if \(i\) is even.

  In every case \(q_{B,p}\) and \(q_{Z,p}\) are isomorphic if and only if \(i\)
  is even, which is exactly the condition of \cref{prop:padic-reduction} for
  \(q_Z\) to be positive definite. This proves the theorem.
\end{proof}

\begin{remark}[Hypotheses of the main theorem]
  \label{rem:mainthm-hypotheses}
  \uses{def:mainthm-setup, def:exotic-motive}
  % SOURCE: [Ancona], §8, Remark (hypothesismainthm) (read from references/source/fourfolds_2020.tex)
  % SOURCE QUOTE: "\begin{enumerate}
  % \item The main example of such a motive $M$ we have in mind is an exotic motive of an abelian fourfold (see Section \ref{sectionexotic}). Another example coming from geometry is given in Proposition \ref{propfermat}.
  % \item\label{homvsrat} One can actually make the  hypothesis a little bit more flexible and work with homological motives instead of Chow motives. This will not matter for our application to abelian varieties.
  % \end{enumerate}"
  \textit{Source: Ancona, arXiv:1806.03216, §8.}
  The main example of such a motive \(M\) is an exotic motive of an abelian
  fourfold (see \cref{def:exotic-motive}). One may relax the hypothesis slightly
  and work with homological motives instead of Chow motives; this does not
  matter for the application to abelian varieties.
\end{remark}

\begin{proposition}[Consequences of the main hypotheses]
  \label{prop:mainthm-properties}
  \lean{MR4199442StandardConjecturesForAbelianFourfolds.mainTheoremProperties}
  \uses{def:mainthm-setup, def:realizations, def:num-motives, prop:motive-num-bound}
  % SOURCE: [Ancona], §8, Proposition (tortura) (read from references/source/fourfolds_2020.tex)
  % SOURCE QUOTE: "Under the hypothesis of Theorem \ref{mainthm} the following holds:
  % \begin{enumerate}
  % \item The quadratic form $q_Z$   on $V_Z$ is non-degenerate.
  % \item  Given a classical realization $R$, the vector space $R(M_{\vert{k}})$  is of dimension two.
  % \item Given a classical realization $R$, the vector space $R(M_{\vert{k}})$ is spanned by algebraic classes.
  % \item Numerical equivalence on $\Hom_{\CHM(k)_\Q}(\one,M_{\vert{k}})$ coincides with the homological equivalence for any classical cohomology.
  % \item\label{giorno} Fix a classical realization $R$ and call $L$ the field of coefficients of  $R$. Then the equality $V_Z\otimes _\Q L = R(M)$ holds.
  % \end{enumerate}"
  \textit{Source: Ancona, arXiv:1806.03216, §8.}
  Under the hypotheses of \cref{thm:mainthm} the following hold.
  \begin{enumerate}
    \item The quadratic form \(q_Z\) on \(V_Z\) is non-degenerate.
    \item For every classical realization \(R\), the vector space
      \(R(M_{\vert{k}})\) has dimension two.
    \item For every classical realization \(R\), the vector space
      \(R(M_{\vert{k}})\) is spanned by algebraic classes.
    \item Numerical equivalence on \(\Hom_{\CHM(k)_\Q}(\one, M_{\vert{k}})\)
      coincides with homological equivalence for any classical cohomology.
    \item If \(R\) is a classical realization with coefficient field \(L\), then
      \(V_Z \otimes_\Q L = R(M)\).
  \end{enumerate}
\end{proposition}

% SOURCE QUOTE PROOF: "Let us start with (2). By the comparison theorem between singular and $\ell$-adic cohomology, we have ${\dim_\Q V_B = \dim_{\Q_\ell}R_\ell(M_{\vert{\C}})}$, for all primes $\ell$, including $\ell=p$.  Then, by smooth proper base change we have   $\dim_{\Q_\ell}R_\ell(M_{\vert{\C}})=\dim_{\Q_\ell}R_\ell(M_{\vert{k}})$ for all $\ell\neq p$ and finally by the $p$-adic comparison theorem we have
%  $\dim_{\Q_p}R_p(M_{\vert{\C}})=\dim_{\Frac(W(k))}R_\crys(M_{\vert{k}})$. This concludes (2) as, by hypothesis, $\dim_\Q V_B=2$.
%  The proof of (3)-(5) goes as in Proposition \ref{motivisupersingolari}. There, the hypothesis that the motive was of abelian type was used to ensure that all  realizations have the same dimension (see Remark \ref{remarkrank}). Here, this is replaced by part (2).
%  Let us now show part (1). First notice that it is enough to show that $R(q_{\vert{k}})$ is non-degenerate (for some classical realization) because of parts (2)-(5).
%  Then, again by using the comparison theorems, this is equivalent to the fact that $R_B(q)$ is non-degenerate. On the other hand this is the case as $R_B(q)$ is a polarization of Hodge structures."
\begin{proof}
  \uses{def:realizations, prop:motive-num-bound, def:mainthm-setup}
  We start with (2). By the comparison theorem between singular and
  \(\ell\)-adic cohomology, \(\dim_\Q V_B = \dim_{\Q_\ell} R_\ell(M_{\vert\C})\)
  for all primes \(\ell\), including \(\ell = p\). By smooth proper base change,
  \(\dim_{\Q_\ell} R_\ell(M_{\vert\C}) = \dim_{\Q_\ell} R_\ell(M_{\vert{k}})\)
  for all \(\ell \neq p\), and by the \(p\)-adic comparison theorem
  \(\dim_{\Q_p} R_p(M_{\vert\C}) = \dim_{\Frac(W(k))} R_\crys(M_{\vert{k}})\).
  Since \(\dim_\Q V_B = 2\) by hypothesis, this gives (2).

  Parts (3)--(5) follow exactly as in \cref{prop:motive-num-bound}: the
  hypothesis used there that the motive is of abelian type, ensuring that all
  realizations have the same dimension, is here replaced by part (2).

  Finally, for part (1) it is enough to show that \(R(q_{\vert{k}})\) is
  non-degenerate for some classical realization, by parts (2)--(5). Using the
  comparison theorems once more, this is equivalent to \(R_B(q)\) being
  non-degenerate, which holds because \(R_B(q)\) is a polarization of Hodge
  structures.
\end{proof}

\section{Quadratic forms}

\subsection*{Rank-two quadratic forms}

We recall the classical facts on quadratic forms that reduce \cref{thm:mainthm}
to a \(p\)-adic question (\cref{prop:padic-reduction}). We work only in the
context needed later: non-degenerate \(\Q\)-quadratic forms of rank \(2\).

\begin{definition}[Hilbert symbol]
  \label{def:hilbert-symbol}
  \lean{MR4199442StandardConjecturesForAbelianFourfolds.hilbertSymbol}
  % SOURCE: [Ancona], §9, Definition (Hilbert) (read from references/source/fourfolds_2020.tex)
  % SOURCE QUOTE: "Let $q$ be a $\Q$-quadratic form of rank $2$. Define $\varepsilon_\nu(q)$, the Hilbert symbol of $q$ at $\nu$, as  $+1$ if the equation
  % \[x^2-q(y,z)=0\]
  % has a nonzero solution in $x,y,z\in \Q_\nu$, and as $-1$ otherwise. Depending on the context we may write $\varepsilon_p(q)$ or $\varepsilon_\R(q)$."
  \textit{Source: Ancona, arXiv:1806.03216, §9.}
  Let \(q\) be a \(\Q\)-quadratic form of rank \(2\) and \(\nu\) a place of
  \(\Q\). The \emph{Hilbert symbol} \(\varepsilon_\nu(q)\) of \(q\) at \(\nu\) is
  \(+1\) if the equation
  \[x^2 - q(y,z) = 0\]
  has a nonzero solution \(x, y, z \in \Q_\nu\), and \(-1\) otherwise. According
  to context we write \(\varepsilon_p(q)\) or \(\varepsilon_\R(q)\).
\end{definition}

\begin{remark}[Signature of a rank-two form]
  \label{rem:rank2-signature}
  \uses{def:hilbert-symbol}
  % SOURCE: [Ancona], §9, Remark (remarksignature) (read from references/source/fourfolds_2020.tex)
  % SOURCE QUOTE: "Let $q$ be a $\Q$-quadratic form of rank $2$. It is positive definite if and only if its discriminant is positive and $\varepsilon_\R(q)=+1$
  % and it is negative definite if and only if its discriminant is positive and $\varepsilon_\R(q)=-1.$"
  \textit{Source: Ancona, arXiv:1806.03216, §9.}
  A \(\Q\)-quadratic form \(q\) of rank \(2\) is positive definite if and only
  if its discriminant is positive and \(\varepsilon_\R(q) = +1\), and it is
  negative definite if and only if its discriminant is positive and
  \(\varepsilon_\R(q) = -1\).
\end{remark}

\begin{proposition}[Local invariants determine a rank-two form over $\Q_p$]
  \label{prop:local-invariant}
  \lean{MR4199442StandardConjecturesForAbelianFourfolds.rank2FormIsoQp}
  \uses{def:hilbert-symbol}
  % SOURCE: [Ancona], §9, Proposition (localinvariant), citing [Serre, Cours d'arithmétique, §2.3] (read from references/source/fourfolds_2020.tex)
  % SOURCE QUOTE: "Let $p$ be a prime number and $q_1$ and $q_2$   two non-degenerate $\Q$-quadratic forms of rank $2$. Then
  % \[q_1 \otimes \Q_p \cong q_2 \otimes \Q_p\]
  % if and only if the discriminants of $q_1$ and $q_2$ coincide in $\Q_p^*/(\Q_p^*)^2$ and
  % \[\varepsilon_p(q_1)=\varepsilon_p(q_2).\]"
  \textit{Source: Ancona, arXiv:1806.03216, §9 (Serre, \emph{Cours d'arithmétique}, §2.3).}
  Let \(p\) be a prime and \(q_1, q_2\) two non-degenerate \(\Q\)-quadratic
  forms of rank \(2\). Then
  \[q_1 \otimes \Q_p \cong q_2 \otimes \Q_p\]
  if and only if the discriminants of \(q_1\) and \(q_2\) coincide in
  \(\Q_p^*/(\Q_p^*)^2\) and \(\varepsilon_p(q_1) = \varepsilon_p(q_2)\).
\end{proposition}

\begin{proof}
  \uses{def:hilbert-symbol}
  \textit{Source: Serre, Cours d'arithmétique, §2.3 (classical; proof sketch
  supplied).}
  Over a \(p\)-adic field a non-degenerate quadratic form of a given rank is
  classified up to isomorphism by two invariants: its discriminant, viewed in
  \(\Q_p^*/(\Q_p^*)^2\), and its Hasse--Witt invariant. Two such forms of equal
  rank over \(\Q_p\) are therefore isomorphic if and only if both invariants
  agree. For a rank-two form the Hasse--Witt invariant reduces to a single
  Hilbert symbol, which — by the defining equation \(x^2 - q(y,z) = 0\) of
  \cref{def:hilbert-symbol} — is exactly \(\varepsilon_p(q)\). Hence for
  rank-two forms the pair (discriminant in \(\Q_p^*/(\Q_p^*)^2\),
  \(\varepsilon_p\)) is a complete system of invariants, and
  \(q_1 \otimes \Q_p \cong q_2 \otimes \Q_p\) if and only if the discriminants
  of \(q_1\) and \(q_2\) coincide in \(\Q_p^*/(\Q_p^*)^2\) and
  \(\varepsilon_p(q_1) = \varepsilon_p(q_2)\).
\end{proof}

\begin{theorem}[Product formula for the Hilbert symbol]
  \label{thm:product-formula}
  \lean{MR4199442StandardConjecturesForAbelianFourfolds.hilbertProductFormula}
  \uses{def:hilbert-symbol}
  % SOURCE: [Ancona], §9, Theorem (productformula), citing [Serre, Cours d'arithmétique, §3.1] (read from references/source/fourfolds_2020.tex)
  % SOURCE QUOTE: "Let $q$ be non-degenerate $\Q$-quadratic form of rank $2$ and $\varepsilon_\nu(q)$ be as in Definition \ref{Hilbert}. Then for all but finite places $\nu$ the equality $\varepsilon_\nu(q)=+1$ holds. Moreover, the following product formula running on all places holds
  % \[\prod_\nu  \varepsilon_\nu(q)= +1.\]"
  \textit{Source: Ancona, arXiv:1806.03216, §9 (Serre, \emph{Cours d'arithmétique}, §3.1).}
  Let \(q\) be a non-degenerate \(\Q\)-quadratic form of rank \(2\), with Hilbert
  symbols \(\varepsilon_\nu(q)\) as in \cref{def:hilbert-symbol}. Then
  \(\varepsilon_\nu(q) = +1\) for all but finitely many places \(\nu\), and the
  product formula over all places holds:
  \[\prod_\nu \varepsilon_\nu(q) = +1.\]
\end{theorem}

\begin{proof}
  \uses{def:hilbert-symbol}
  \textit{Source: Serre, Cours d'arithmétique, §3.1 (classical; proof sketch
  supplied).}
  Up to scaling, diagonalize \(q = \langle a, b \rangle\) with
  \(a, b \in \Q^*\). By \cref{def:hilbert-symbol} the local invariant
  \(\varepsilon_\nu(q)\) is then the Hilbert symbol \((a,b)_\nu\) of \(\Q\) at
  the place \(\nu\), measuring whether \(z^2 = a x^2 + b y^2\) has a nonzero
  solution over \(\Q_\nu\). For a place \(\nu\) not lying over \(2\), \(\infty\),
  or a prime dividing \(ab\), both \(a\) and \(b\) are units in \(\Z_\nu\) and
  the form \(\langle a, b \rangle\) represents every unit, so
  \((a,b)_\nu = +1\); this leaves only the finitely many places dividing
  \(2 a b \infty\), giving the finiteness assertion. The global identity
  \[\prod_\nu \varepsilon_\nu(q) = \prod_\nu (a,b)_\nu = +1\]
  is Hilbert reciprocity for \(\Q\), which is equivalent to the law of quadratic
  reciprocity together with its supplementary laws.
\end{proof}

\begin{corollary}[Sign comparison from $\ell$-adic data]
  \label{cor:signature-top}
  \lean{MR4199442StandardConjecturesForAbelianFourfolds.signatureComparison}
  \uses{def:hilbert-symbol, rem:rank2-signature, prop:local-invariant, thm:product-formula}
  % SOURCE: [Ancona], §9, Corollary (signaturetop) (read from references/source/fourfolds_2020.tex)
  % SOURCE QUOTE: "Let $q_1$ and $q_2$  be two non-degenerate $\Q$-quadratic forms of rank $2$ and let $p$ be a prime number. Suppose that, for all primes $\ell$ different from $p$, we have
  % \[q_1 \otimes \Q_\ell \cong q_2 \otimes \Q_\ell.\]
  % Then $q_1$ is positive definite if and only if one of the following two cases happens:
  % \begin{enumerate}
  % \item  The quadratic forms $q_1\otimes \Q_p$ and $q_2\otimes \Q_p$ are isomorphic and $q_2$ is positive definite.
  % \item  The quadratic forms $q_1\otimes \Q_p$ and $q_2\otimes \Q_p$ are not isomorphic and $q_2$ is negative definite.
  % \end{enumerate}"
  \textit{Source: Ancona, arXiv:1806.03216, §9.}
  Let \(q_1, q_2\) be two non-degenerate \(\Q\)-quadratic forms of rank \(2\),
  and let \(p\) be a prime. Suppose that
  \[q_1 \otimes \Q_\ell \cong q_2 \otimes \Q_\ell \qquad
    \text{for all primes } \ell \neq p.\]
  Then \(q_1\) is positive definite if and only if one of the following holds:
  \begin{enumerate}
    \item \(q_1 \otimes \Q_p\) and \(q_2 \otimes \Q_p\) are isomorphic and
      \(q_2\) is positive definite; or
    \item \(q_1 \otimes \Q_p\) and \(q_2 \otimes \Q_p\) are not isomorphic and
      \(q_2\) is negative definite.
  \end{enumerate}
\end{corollary}

% SOURCE QUOTE PROOF: "The $\ell$-adic hypothesis implies in particular that the discriminants of $q_1$ and $q_2$ coincide in $ \Q_\ell^*/(\Q_\ell^*)^2$
%  for all   $\ell\neq p$. This implies that they coincide in $ \Q^*/(\Q^*)^2$
% by \cite[Theorem 3 in 5.2]{Ireland}.
% If the discriminants are negative, none of the conditions in the statement holds and the equivalence is clear. From now on we suppose that the discriminants are   positive. By Remark \ref{remarksignature}, $q_1$ is positive definite if and only $\varepsilon_\R(q_1)=+1.$
% There are then two cases, namely $q_2$ is positive definite, respectively $q_2$ is negative definite. Again by Remark \ref{remarksignature} they are equivalent to $\varepsilon_\R(q_2)=+1$, respectively $\varepsilon_\R(q_2)=-1$.
% Now, Theorem \ref{productformula} implies that
% $\prod_\nu  \varepsilon_\nu(q_1)=\prod_\nu  \varepsilon_\nu(q_2).$
% Combining this with the $\ell$-adic isomorphisms we deduce
% \[\varepsilon_\R(q_1)\varepsilon_p(q_1)=\varepsilon_\R(q_2)\varepsilon_p(q_2).\] This means that $q_1$ is positive definite if and only if
%  \[\varepsilon_p(q_1)=\varepsilon_\R(q_2)\varepsilon_p(q_2).\]
%  This relation is equivalent to the fact that one of the following two situations hold:
%  \begin{enumerate}
% \item   $q_2$ is positive definite and $\varepsilon_p(q_1)=\varepsilon_p(q_2).$
% \item  $q_2$ is negative definite and $\varepsilon_p(q_1)\neq\varepsilon_p(q_2).$
% \end{enumerate}
% As the discriminant of $q_1$ and $q_2$ coincide, the equality $\varepsilon_p(q_1)=\varepsilon_p(q_2)$ is equivalent to the fact that $q_1\otimes \Q_p$ and $q_2\otimes \Q_p$ are isomorphic (Proposition \ref{localinvariant})."
\begin{proof}
  \uses{rem:rank2-signature, prop:local-invariant, thm:product-formula}
  The \(\ell\)-adic hypothesis implies that the discriminants of \(q_1\) and
  \(q_2\) coincide in \(\Q_\ell^*/(\Q_\ell^*)^2\) for all \(\ell \neq p\), and
  hence (by the Hasse principle for squares in \(\Q\)) coincide already in
  \(\Q^*/(\Q^*)^2\).

  If the common discriminant is negative, none of the listed conditions holds
  and neither does positive definiteness, so the equivalence is clear. Assume
  henceforth the discriminant is positive. By \cref{rem:rank2-signature},
  \(q_1\) is positive definite if and only if \(\varepsilon_\R(q_1) = +1\), and
  \(q_2\) is positive (resp.\ negative) definite if and only if
  \(\varepsilon_\R(q_2) = +1\) (resp.\ \(-1\)).

  By \cref{thm:product-formula},
  \(\prod_\nu \varepsilon_\nu(q_1) = \prod_\nu \varepsilon_\nu(q_2) = +1\).
  Combining this with the \(\ell\)-adic isomorphisms (which give
  \(\varepsilon_\ell(q_1) = \varepsilon_\ell(q_2)\) for all \(\ell \neq p\))
  leaves only the real and \(p\)-adic places, whence
  \[\varepsilon_\R(q_1)\,\varepsilon_p(q_1)
    = \varepsilon_\R(q_2)\,\varepsilon_p(q_2).\]
  Thus \(q_1\) is positive definite, i.e.\ \(\varepsilon_\R(q_1) = +1\), if and
  only if \(\varepsilon_p(q_1) = \varepsilon_\R(q_2)\,\varepsilon_p(q_2)\); that
  is, if and only if either \(q_2\) is positive definite and
  \(\varepsilon_p(q_1) = \varepsilon_p(q_2)\), or \(q_2\) is negative definite
  and \(\varepsilon_p(q_1) \neq \varepsilon_p(q_2)\). Since the discriminants of
  \(q_1\) and \(q_2\) coincide, the equality
  \(\varepsilon_p(q_1) = \varepsilon_p(q_2)\) is, by
  \cref{prop:local-invariant}, equivalent to \(q_1 \otimes \Q_p\) and
  \(q_2 \otimes \Q_p\) being isomorphic. This is the assertion.
\end{proof}

\subsection*{Reduction to a $p$-adic comparison}

\begin{proposition}[$p$-adic reduction of the main theorem]
  \label{prop:padic-reduction}
  \lean{MR4199442StandardConjecturesForAbelianFourfolds.padicReduction}
  \uses{def:mainthm-setup, prop:mainthm-properties, cor:signature-top}
  % SOURCE: [Ancona], §9, Proposition (padicreduction) (read from references/source/fourfolds_2020.tex)
  % SOURCE QUOTE: "Let us keep notation from Theorem \ref{mainthm}. Let $p$ be the characteristic of $k$ and $i$ be the unique non-negative integer such that the Hodge structure $V_B$ is of type $(i,-i)$ and  $(-i,i)$. Then, the quadratic form $q_Z$ is positive definite if and only if the  following   holds:
  % \begin{enumerate}
  % \item When $i$ is even: the quadratic forms $q_B\otimes \Q_p$ and $q_Z\otimes \Q_p$ are isomorphic.
  % \item  When $i$ is odd: the quadratic forms $q_B\otimes \Q_p$ and $q_Z\otimes \Q_p$ are not isomorphic.
  % \end{enumerate}"
  \textit{Source: Ancona, arXiv:1806.03216, §9.}
  Keep the notation of \cref{def:mainthm-setup}. Let \(p\) be the characteristic
  of \(k\), and let \(i\) be the unique non-negative integer such that the Hodge
  structure \(V_B\) is of type \((i,-i)\) and \((-i,i)\). Then \(q_Z\) is
  positive definite if and only if:
  \begin{enumerate}
    \item when \(i\) is even, the quadratic forms \(q_B \otimes \Q_p\) and
      \(q_Z \otimes \Q_p\) are isomorphic; and
    \item when \(i\) is odd, the quadratic forms \(q_B \otimes \Q_p\) and
      \(q_Z \otimes \Q_p\) are not isomorphic.
  \end{enumerate}
\end{proposition}

% SOURCE QUOTE PROOF: "By  Proposition \ref{tortura}(\ref{giorno}), we have that $q_Z \otimes \Q_\ell = R_\ell(q_{\vert{k}})$. By the comparison theorem, we have $q_B \otimes \Q_\ell = R_\ell(q_{\vert{\C}})$. Combining these equalities with smooth proper base change in $\ell$-adic cohomology   we deduce that \[q_B \otimes \Q_\ell = q_Z \otimes \Q_\ell.\]
% On the other hand, by hypothesis, $q_B$ is a polarization for the Hodge structure $V_B$ hence it is positive definite if $i$ is even and negative definite if $i$ is odd. We can now conclude by applying Corollary \ref{signaturetop} to $q_1=q_Z$ and $q_2=q_B$."
\begin{proof}
  \uses{prop:mainthm-properties, cor:signature-top}
  By \cref{prop:mainthm-properties}\,(part~(5)) we
  have \(q_Z \otimes \Q_\ell = R_\ell(q_{\vert{k}})\), and by the comparison
  theorem \(q_B \otimes \Q_\ell = R_\ell(q_{\vert\C})\). Combining these with
  smooth proper base change in \(\ell\)-adic cohomology yields
  \[q_B \otimes \Q_\ell = q_Z \otimes \Q_\ell \qquad
    \text{for all } \ell \neq p.\]
  On the other hand, by hypothesis \(q_B\) is a polarization of the Hodge
  structure \(V_B\), hence \(q_B\) is positive definite when \(i\) is even and
  negative definite when \(i\) is odd. The claim now follows by applying
  \cref{cor:signature-top} with \(q_1 = q_Z\) and \(q_2 = q_B\).
\end{proof}
